% !TeX program = xelatex
% !Mode:: "TeX:UTF-8"
%% ═══════════════════════════════════════════════════════════
%%  全国大学生数学建模竞赛 (CUMCM) LaTeX 模板
%%  基于 LaTeX工作室 数维杯模板 v4 (2025/04/26) 改编
%%  编译方式: xelatex main.tex (运行两次以更新交叉引用)
%% ═══════════════════════════════════════════════════════════
%
% ========== 排版规则速查 ==========
% 1. 编译：xelatex × 3 次（第一次写aux，第二/三次解析）
% 2. 浮动控制：\usepackage{placeins} + 每\subsubsection前\FloatBarrier
% 3. 避免大段空白：图表用[htbp]而非[H]（符号表除外）
% 4. 公式编号：全文自动连续编号，手写"式(X)"需与自动编号一致
% 5. 算法伪代码：三线表环境，\footnotesize字号，p{0.75\textwidth}
% 6. 摘要：独立一页，首行缩进2em，三类加粗（模型名+方法名+核心数值）
% 7. 图表：图名在下，表名在上，每图后≥100字分析
% =================================

\documentclass[a4paper]{cumcmthesis}

% ═══ 基础宏包 ═══
\usepackage{graphicx}
\usepackage{placeins}       % 提供\FloatBarrier（防图表跨子节）
\usepackage{booktabs, colortbl}
\usepackage{longtable, multirow, tabularx}
\usepackage{float}          % 提供[H]（仅符号表使用）
\usepackage{xcolor}
\usepackage{tikz}
\usepackage{pgfplots}
\pgfplotsset{compat=1.18}
\usetikzlibrary{shapes.geometric, arrows.meta, positioning, calc}

% ═══ 竞赛配置 ═══
\mcmsetup{
    CTeX = true,
    tcn  = {\xiaowuhao 2026A00001},  % ← 替换为实际队号
    problem = A,                       % ← 替换为实际题号
    sheet = true,
    titleinsheet = false,
    keywordsinsheet = false,
    titlepage = true,
}

% ═══ 自定义命令 ═══
% 使用说明：摘要中的加粗用\textbf{}，正文中的加粗仅保留给摘要
\newcommand{\E}{\mathbb{E}}
\newcommand{\cC}{\mathcal{C}}
\newcommand{\cO}{\mathcal{O}}
\newcommand{\cS}{\mathcal{S}}
\newcommand{\cI}{\mathcal{I}}
\newcommand{\cT}{\mathcal{T}}
\newcommand{\cM}{\mathcal{M}}
\newcommand{\ftilde}{\tilde{f}}

% ============================================================
%  正文开始
% ============================================================
\begin{document}

% ═══════════════ 1. 摘要 ═══════════════
% 【规则】
% - 独立一页（国赛第3页，前2页为承诺书+编号页）
% - 首行缩进2em（\setlength{\parindent}{2em}）
% - 首段：只写"问题+模型+方法+目标"四要素
% - 每子问题段：3-5行，含模型名（加粗）+方法+数值（加粗）
% - 摘要页顶部为论文标题（三号黑体居中）
% - 关键词加粗，用分号分隔
% - 最后一句必须是总结性结论，不以灵敏度结尾
\setlength{\parindent}{2em}

\begin{abstract}
摘要在各段中加粗以下三类：
（1）自创指标/模型名：货位热度指标、货格配对策略、二分图最大权匹配
（2）核心方法：遗传算法、Hungarian算法
（3）关键数值：20.98s、99.0%、4.78h、1640对

注意：正文段落中不加粗任何概念名，加粗只在摘要中出现。
\end{abstract}

% ═══════════════ 2. 问题重述 ═══════════════
% 【规则】
% - 不分小节（无\subsection）
% - 连续两段：首段背景+数据参数，次段五问题各一句话
% - 不用\begin{enumerate}列问题
\section{问题重述}

% ═══════════════ 3. 问题分析 ═══════════════
% 【规则】
% - 每题一段（\subsection{问题X的分析}），点明数学本质和求解思路
% - 不写具体数据规模
\section{问题分析}
\subsection{问题一的分析}
\subsection{问题二的分析}
\subsection{问题三的分析}
\subsection{问题四的分析}
\subsection{问题五的分析}

% ═══════════════ 4. 模型假设 ═══════════════
% 【规则】
% - 5条核心假设
% - 直接写结论，不写"为使模型具有可操作性..."等引导句
% - 不写"依据："理由
% - 不写不重要假设（如"系统连续运行""不考虑故障"）
\section{模型假设}
\begin{assumption}
堆垛机水平与垂直运动独立驱动。
\end{assumption}

% ═══════════════ 5. 符号说明 ═══════════════
% 【规则】
% - 三线表
% - 不用float环境（[H]固定），确保紧跟假设之后
\section{符号说明}
\begin{table}[H]
\centering
\caption{符号说明}
\label{tab:symbols}
\begin{tabular}{cll}
\toprule
符号 & 含义 & 单位 \\
\midrule
\end{tabular}
\end{table}

% ═══════════════ 6. 模型的建立与求解 ═══════════════
% 【规则】
% - 8节国赛结构中的核心节（20-22页）
% - 每个问题必须有3个子节：模型建立→算法设计→结果分析
% - 每个子节前加\FloatBarrier防止图表跨子节
% - 子问题间必须有过渡句解释"为什么需要下一个问题"
\section{模型的建立与求解}

\subsection{数据预处理}
% 1页：统计表+箱型分布

% ═══ 问题一 ═══
\subsection{问题一：XX优化}

\FloatBarrier\subsubsection{模型建立}
% 【内容】
% - 决策变量定义
% - 目标函数（编号公式）
% - 约束条件逐条：命名(\emph{约束名}) + 编号公式 + 物理解释（硬/紧/松/软）
% - 可用\textbf{1.}\textbf{2.}\textbf{3.}分子节

\FloatBarrier\subsubsection{算法设计}
% 【格式】
% - 三线表伪代码
% - \begin{table}[htbp] \centering\footnotesize
% - \begin{tabular}{p{0.75\textwidth}}
% - Algorithm N: Title
% - Input:/Output: 行
% - 最后一行含时间复杂度
% - 表后有文字解释（≥1段）

\FloatBarrier\subsubsection{求解结果与分析}
% 【内容】
% - 数值结果表（三线表[htbp]）
% - 图表（≥1张，[htbp]浮动）
% - ≥100字图分析：展示什么+关键发现+物理/工程含义
% - 与前问对比

% ═══ 问题二-五（结构同上） ═══
% 每两个问题之间必须有过渡段。禁用"节末核心发现"和"综上"。

% ═══════════════ 7. 灵敏度分析 ═══════════════
% 【规则】
% - 含单参数扰动表+联合扰动+机理分析
% - 模型检验作为子节
\section{灵敏度分析}

% ═══════════════ 8. 模型评价与推广 ═══════════════
% 【规则】
% - 含：问题递进关系→GA差异分析→线性速度适用范围→约束可行性验证→优缺点→推广
% - 劣势要与假设呼应
% - 禁用："模型有优点也有缺点"、"本文模型还比较粗糙"
\section{模型评价与推广}

% ═══════════════ 9. 结论 ═══════════════
% 【规则】
% - ≥5条编号总结（\begin{enumerate}）
% - 含数值汇总表
% - 最后一句概括全部工作
\section{结论}

% ═══════════════ 10. 参考文献 ═══════════════
% 【规则】
% - ≥10篇，格式统一
\bibliographystyle{...}
\bibliography{ref.bib}

% ═══════════════ 11. 附录 ═══════════════
% 【规则】
% - 每问核心代码各一个listing
% - 代码必须与正文算法一致
\begin{appendices}
\textbf{程序一：}...
\end{appendices}

\end{document}
